%! Semi-log Plots — Unit 5, Lesson 5.8 — Guided Notes
\documentclass[10pt]{article}
\usepackage{precalculus-article}
\usepackage{precalculus-boxes}
\newcommand{\TallMath}[1]{$\displaystyle #1\rule[-1.4em]{0pt}{3.2em}$}

\begin{document}

\pageheader{Unit 5, Lesson 5.8}{Guided Notes}
\namedateperiod

\vspace{0.1in}

\begin{objectivebox}
\textbf{LO 2.15.A:} Determine if an exponential model is appropriate by examining a
semi-log plot. \quad \textbf{LO 2.15.B:} Construct the linearization of exponential data.

By the end of this lesson I will be able to:
\begin{itemize}[leftmargin=1.5em,itemsep=3pt,topsep=3pt]
  \item Explain what a semi-log plot is and identify when exponential data appears linear
        on one (EK 2.15.A.1).
  \item Write the linearized form $\log_n(y) = (\log_n b)\cdot x + \log_n a$ for a model
        $y = ab^x$ (EK 2.15.B.2).
  \item Recover $a$ and $b$ from the slope and intercept of the linearized data (EK 2.15.B.1).
\end{itemize}
\end{objectivebox}

\vspace{0.08in}

\begin{vocabbox}
\termblanklong{semi-log plot}
\termblanklong{logarithmic scale}
\termblanklong{linearization}
\termblanklong{slope of linearized data}
\termblanklong{$y$-intercept of linearized data}
\end{vocabbox}

\vspace{0.08in}

\begin{hookbox}
\textbf{Consider:} A scientist plots bacterial population data on a regular graph and on a
semi-log graph.  On the regular graph the curve bends sharply upward.  On the semi-log graph
the same data forms a straight line.

\vspace{4pt}
What does this tell the scientist about the type of growth model? \writeline
\end{hookbox}

\vspace{0.08in}

\begin{notesbox}{What Is a Semi-log Plot? (EK 2.15.A.1--2)}
\textbf{Definition:} A semi-log plot has one axis that is \blank{4cm} scaled.  On a
$y$-axis semi-log plot, the tick marks represent $10^0, 10^1, 10^2, \ldots$ (equal spacing
= equal \emph{ratios}).

\vspace{6pt}
\textbf{Key fact:} Data or functions with exponential characteristics appear
\blank{3cm} on a semi-log plot.

\vspace{6pt}
\textbf{Advantage (EK 2.15.A.2):} A constant \blank{3cm} needs to be added to
the dependent variable values to reveal that an exponential model is appropriate.
(Compare to the Lesson 2.6 approach.)
\end{notesbox}

\vspace{0.08in}

\begin{notesbox}{The Linearization of $y = ab^x$ (EK 2.15.B.2)}
Start with $y = ab^x$.  Apply $\log_{10}$ to both sides:
\[
  \log_{10}(y) = \log_{10}(a \cdot b^x) = \log_{10}(a) + \log_{10}(b^x) = \log_{10}(a) + x\,\log_{10}(b).
\]

So the linearized model is:
\[
  \log_{10}(y) \;=\; \bigl(\log_{10} b\bigr)\cdot x \;+\; \log_{10} a.
\]
\begin{itemize}[leftmargin=1.5em,itemsep=4pt,topsep=4pt]
  \item Slope $=$ \blank{3.5cm}
  \item $y$-intercept $=$ \blank{3.5cm}
\end{itemize}

\vspace{6pt}
\textbf{Example.} For $y = 5 \cdot 3^x$, the linearized model is:

\vspace{4pt}
$\log_{10}(y) = (\log_{10} 3)\cdot x + \log_{10} 5 \approx $ \blank{1.5cm}$\cdot x + $ \blank{1.5cm}
\end{notesbox}

\vspace{0.08in}

\begin{notesbox}{Recovering the Exponential Model (EK 2.15.B.1)}
Given the slope $m$ and $y$-intercept $c$ of the linearized semi-log data:
\[
  b = 10^m \qquad \text{and} \qquad a = 10^c.
\]

\textbf{Example.}  A semi-log plot of $(x,\,\log_{10}(y))$ data has slope $0.602$ and
$y$-intercept $0.699$.

\vspace{4pt}
$b = 10^{0.602} \approx $ \blank{1.5cm} \qquad $a = 10^{0.699} \approx $ \blank{1.5cm}

\vspace{4pt}
Exponential model: $y = $ \blank{4cm}
\end{notesbox}

\end{document}
